\section{Conclusion}
\label{sec:conclusion}

This paper has shown that interoperability of chemical substance data in
European regulation is fundamentally limited by the absence of a shared and
formalised notion of chemical identity. Across 18 regulatory datasets, a
substantial fraction of substances cannot be linked to a defined molecular
structure, and commonly used identifiers such as CAS numbers do not provide a
reliable basis for cross-dataset integration.

We demonstrate that this heterogeneity leads to unresolved semantic relations
between regulatory entities and creates a combinatorial barrier to
retrospective harmonisation. Data-driven approaches, including semantic
embeddings, are insufficient to recover chemical identity in the absence of
structure-based representations.

At the same time, our results show that interoperability becomes feasible when
structure-based identifiers such as InChIKey are combined with ontology-driven
approaches and knowledge graph representations. This highlights that the
interoperability problem is not primarily technical, but originates in how
chemical substances are defined and represented in regulatory frameworks.

We therefore argue that achieving FAIR and interoperable chemical data in
Europe requires a shift towards structure-based identifiers, formal
ontologies, and linked data principles embedded directly in regulatory data
models and legislation. Without such changes, the gap between regulatory data
and scientific data infrastructures will persist, limiting the effectiveness of
data-driven environmental policy.

Future work should focus on developing hybrid semantic models for mixtures and
substance groups, as well as on translating these principles into operational
standards for regulatory data publication.